\section{\uppercase{Related Work}}
\label{sec:related}
\subsection{Distributional and Rank-Based Coding}
Neural linguistic steganography uses an autoregressive distribution as an information-bearing alphabet. Neural Linguistic Steganography \cite{ziegler} encodes a message through arithmetic decoding and recovers it by replaying arithmetic encoding. That work connects information efficiency to the model's conditional distributions. Our comparator follows that approach with an explicitly identified finite-message adaptation.\aifootnote

Distribution preservation is a stronger objective than recovering a packet or producing fluent text. The coupling formulation \cite{coupling} characterizes perfectly secure procedures through couplings and evaluates minimum entropy coupling on language, audio and Image Transformer channels. More recent work \cite{range} proposes rotation-based range coding with distribution-preservation guarantees and high entropy utilization. These distribution-preserving methods were not experimentally matched here. Our results establish neither superiority to them nor a limitation of every arithmetic or range coder.

Calgacus \cite{calgacus} instead transfers each payload token's context-dependent rank to a cover context. It discusses separate models, different vocabulary sizes and other discrete domains. RankCloak \cite{rankcloak} adapts rank transport to serialized cryptographic artifacts and implements bounded byte coding and entropy gating. Its repository manuscript reports 6,480 recoveries with preserved token sequences, but 88/144 recoveries in a rendered-text robustness subset. These are prior reported results, not reruns or matched baselines for our study. LlmStenoExplore \cite{explore} reproduces Calgacus and investigates finite-key collisions and perturbation sensitivity. We reuse focused rank, replay and numerical practices from this line of work, while making saved-artifact recovery the primary endpoint and separating the packet key from model conditioning.

\subsection{Pixels and Tokenization}
Pixel-Stega \cite{pixelstega} is the closest image precedent. It obtains autoregressive pixel probabilities from PixelCNN++, applies finite-precision arithmetic stegosampling, and describes transmission through a lossless public channel. Its color-image formulation identifies the red channel as the embedded variable. Our contribution is not the invention of autoregressive image steganography. We implement all delivered R, G and B values as successive observable symbols, with the correct conditional mixture posteriors, and compare three coders using the same authenticated packet and cropped, shared-row canvas as in the reverse text direction. Published Pixel-Stega capacity and detector results use different allocations and cannot be ranked against our net authenticated rates.

Published stepwise verification \cite{consistency} directly addresses detokenization and retokenization inconsistency through stepwise candidate verification. We adapt this existing idea to a pinned byte-exact tokenizer and a top-256 candidate pool. Our development comparison tests the operational consequence of replacing singleton eligibility with complete-prefix consistency, while holding the other rules fixed. It adds paired evidence for this transport contract rather than claiming tokenization consistency as a new technique.

\subsection{Detectability and Execution}
Probability and rank statistics have a history in generated-text analysis. GLTR \cite{gltr}, for example, exposes model-based token statistics to support detection. Our observer does not classify human versus generated text. It distinguishes steganographic carriers from ordinary samples of the same eligible model distribution, using two frozen whole-carrier scores. The conditioning mismatch experiment further asks how those scores change when the exact image model remains known but the observer supplies a different row. It does not test an arbitrary or fully blind adversary.

CUDA graphs are established execution infrastructure \cite{pytorchgraphs}. We measure their benefit for repeated PixelCNN++ replay and verify exact coder-interface and saved-file equivalence. A faster but numerically changed distribution could reorder ranks and break recovery.


